\documentclass[12pt]{article}
\usepackage[utf8]{inputenc}
\usepackage[T1]{fontenc}
\usepackage{hyperref}
\title{Aikan\~a Language Resources for K-12 Curriculum Development}
\author{Compiled by Claude Sonnet 5 (claude-sonnet-5) for Andr\'as Kornai}
\date{}
\begin{document}
\maketitle

\section{Introduction}
Aikan\~a (ISO 639-3: tba; Glottocode: aika1237) is an endangered language isolate of Rond\^onia, Brazil, with SOV word order. Estimates range from about 170 to about 250 speakers. The search found only a few scholarly and profile sources, and no school materials.

\section{Primary Educational Resources}
No primers, textbooks or school programmes for Aikan\~a were found in search results.

\section{Linguistic and Reference Materials}
\begin{itemize}
\item Wikipedia, ``Aikan\~a language'': overview, with a phonological inventory. \url{https://en.wikipedia.org/wiki/Aikan%C3%A3_language}
\item LIAMES (Unicamp) article on Aikan\~a body-part terms (2023). It cites Silva (2012), \emph{Dicion\'ario de ra\'izes da l\'ingua Aikan\~a}, a root dictionary. I did not find the dictionary itself, its publisher, or where to get it. \url{https://periodicos.sbu.unicamp.br/ojs/index.php/liames/article/view/8670923}
\item \emph{Linguagem \& Ensino} (UFPel) article, listed in the search as a study of Aikan\~a oral vowels (Tondineli, 2020). \url{https://periodicos.ufpel.edu.br/index.php/rle/article/download/18288/12004}
\end{itemize}
No complete Aikan\~a grammar was found.

\section{Religious and Community Resources}
\begin{itemize}
\item Povos Ind\'igenas no Brasil (Instituto Socioambiental): Aikan\~a people profile. \url{https://pib.socioambiental.org/en/povo/aikana}
\end{itemize}

\section{Digital and Audio Resources}
\begin{itemize}
\item DoBeS archive (MPI), ``Southeastern Rond\^onia'' documentation project page. The search result did not show its Aikan\~a holdings, so check them before relying on it. \url{https://dobes.mpi.nl/?p=2985}
\end{itemize}

\section{Curriculum Development Context}
The community is small and the language is not known to be taught in schools. A curriculum would need new materials built with speakers, using the root dictionary and the phonological studies above as starting points. Collaboration with researchers at Brazilian universities (such as UNIR) would be sensible.
\end{document}